\subsection{{Mini Services Runner (MSR) \hfill November 2024 -- Present}}
\subtext{\href{https://github.com/bowber/msr}{github.com/bowber/msr}\hfill}
\begin{zitemize}
\item A tool to set up the environment for container-based applications. Built with \textbf{Tauri, SQLite and SolidJS}.
\item Inspired by \textbf{Google Cloud Run}. 
\item This project includes a single desktop app that provides tools to help quickly set up, manage, and maintain an MSR Cluster alongside its services.

\end{zitemize}

\subsection{{MSR Cloud \hfill  Work In Progress}}
\subtext{\href{https://cloud.msrproject.org}{cloud.msrproject.org}\hfill}
\begin{zitemize}
\item This is a work-in-progress commercial project of myself, which leverages MSR to create multiple clusters for customers to deploy their services easily at low cost.
\item It is not yet functional, but it demonstrates my perspective on DevOps and microservices.

\end{zitemize}

\subsection{{axum-db-benchmark \hfill June 2025 -- July 2025}}
\subtext{\href{https://github.com/bowber/axum-db-benchmark/tree/sqlite}{https://github.com/bowber/axum-db-benchmark/tree/sqlite}\hfill}
\begin{zitemize}
\item This is a benchmark of some DB using \textbf{rust + axum + wrk}.
\item It helps me compare each DB for architecture decisions later on.
\item Each DB benchmark code and result is located in its own branch.

\end{zitemize}



\subsection{{Onchain Clash \hfill October 2023 -- June 2024}}
\subtext{\href{https://www.youtube.com/watch?v=oBKoEBgOmE0}{youtube.com/watch?v=oBKoEBgOmE0}\hfill}
\begin{zitemize}
\item A blockchain game built on \textbf{Ancient8}, an ETH gaming Layer 2.
\item This game has been discontinued, but you can still find the marketing video.
\item It can demonstrate my web game development and blockchain skills.

\end{zitemize}

\subsection{{Canvas Image \hfill June 2021 -- August 2024}}
\subtext{\href{https://www.npmjs.com/package/@bowber/canvas-image/v/latest}{npmjs.com/package/@bowber/canvas-image/v/latest}\hfill}
\begin{zitemize}
\item An image interface similar to Python's PIL, based on \texttt{HTMLCanvasElement}.
\item Initially, this was created mainly for migrating \textbf{Rawcutter} from a Python backend to browser-based JavaScript.
\item Later on, I used it for Onchain Clash because I needed a library to generate maps from the pixels of an image.

\end{zitemize}

\subsection{{Rawcutter \hfill March 2020 -- September 2021}}
\subtext{\href{https://rawcutter.com/}{rawcutter.com}\hfill}
\begin{zitemize}
\item A simple image processing tool created for manga translation editors.  
\item This was my first website project, which went through three revisions with three different tech stacks.  
\item The first version was a \textbf{Python Tkinter} desktop app that I built for myself.  
\item The next version was a server-side app using \textbf{Python Flask} with \textbf{vanilla JavaScript}, hosted on \textbf{Heroku}.  
\item The latest version is built with \textbf{React }and hosted on \textbf{GitLab Pages}.  


\end{zitemize}



\subsection{{A Self-host Gitea service.}}
\subtext{\href{https://gitea.msrproject.org}{gitea.msrproject.org}\hfill}
\begin{zitemize}
\item Using \textbf{PostgreSQL} as its database.
\item This demo shows my basic DevOps skills in \textbf{Linux}, \textbf{database hosting}, and \textbf{TLS management}.

\end{zitemize}








